% TODO: Highlights structure 

\begin{frontmatter}

    \begin{abstract}
        
        
        There is an increasing push for operational measures to reduce ships' bunker \ac{fc} and carbon emissions, driven by \ac{imo} mandates. Key performance indicators like \ac{eeoi} focus on fuel efficiency. Trim optimization, virtual arrival, and green routing strategies have emerged. The theoretical basis for all these approaches lies in the accurate predictions of \ac{fc} based on its sailing speed, displacement, trim, climate, and sea state. 

        This study utilized 296 voyage reports from a Bulk Carrier Vessel over one year (from November 16, 2021, to November 21, 2022) and 28 parameters, integrating hydrometeorological big data from \ac{cmems} and \ac{ecmwf} with 19 and 61 parameters, respectively, to evaluate if fusing external public data sources may enhance modeling accuracy as well as highlighting most influential parameters affecting \ac{fc}.
        
        The experimental results revealed that there is a large potentioal of \ac{ml} techniques in predicting ship \ac{fc} accurately by using voyage reports and fusing reliable climate and sea data. However, further validation on similar class of vessels remains necessary to confirm the model's generalizability.
 
    \end{abstract}

    % Research Keywords
    \begin{keyword}
        Voyage Optimization, Bulk Carrier, \ac{fc} Modelling, Machine Learning, CMEMS, ERA5
    \end{keyword}

\end{frontmatter}
